% ID: SEL-D0003-OB-E03
% ESTADO: validada. S004, 2026-11-07; tres pasadas del mismo autor y revisión matemática, didáctica y visual.
\textbf{a) Diferencia e intersección.}\par\nopagebreak[4]
$A-B=A\cap B^c$. Los sucesos $A-B$ y $A\cap B$ son disjuntos y forman $A$:
\[
P(A\cap B)=P(A)-P(A-B)=0.5-0.3=\boxed{0.2}.
\]

\textbf{b) Probabilidad condicionada al complementario.}\par\nopagebreak[4]
Primero obtenemos la probabilidad de $B$ por inclusión-exclusión:
\[
P(B)=P(A\cup B)-P(A)+P(A\cap B)=0.75-0.5+0.2=0.45.
\]
Por tanto, $P(B^c)=0.55>0$ y
\[
P(A/B^c)=P(A\mid B^c)=\frac{P(A\cap B^c)}{P(B^c)}
=\frac{0.3}{0.55}=\boxed{\frac6{11}\approx0.545455}.
\]
El numerador es la diferencia dada, no $P(A\cap B)$.

\textbf{c) Independencia e incompatibilidad.}\par\nopagebreak[4]
No son independientes, porque $P(A)P(B)=0.5(0.45)=0.225\ne0.2=P(A\cap B)$.

Tampoco son incompatibles: su intersección tiene probabilidad 0.2, mientras que dos sucesos incompatibles tendrían intersección vacía y probabilidad 0.

Comprobación de coherencia: las cuatro regiones disjuntas $A\cap B$, $A\cap B^c$, $A^c\cap B$, $A^c\cap B^c$ tienen probabilidades $0.20,0.30,0.25,0.25$. Suman 1 y reconstruyen los datos.
